%!TeX root=../persuasiontop.tex
\chapter[Chapter \thechapter]{}

\begin{figure}[t!]
\centering
\includegraphics[width=\linewidth]{nouveau08}
\end{figure}

\lettrine[lines=4,lraise=0.3]{F}{rom} this time Captain Wentworth and Anne Elliot were repeatedly in the same circle. They were soon dining in company together at Mr Musgrove's, for the little boy's state could no longer supply his aunt with a pretence for absenting herself; and this was but the beginning of other dinings and other meetings.

Whether former feelings were to be renewed must be brought to the proof; former times must undoubtedly be brought to the recollection of each; \textit{they} could not but be reverted to; the year of their engagement could not but be named by him, in the little narratives or descriptions which conversation called forth. His profession qualified him, his disposition lead him, to talk; and »\textit{That} was in the year six;« »\textit{That} happened before I went to sea in the year six,« occurred in the course of the first evening they spent together: and though his voice did not falter, and though she had no reason to suppose his eye wandering towards her while he spoke, Anne felt the utter impossibility, from her knowledge of his mind, that he could be unvisited by remembrance any more than herself. There must be the same immediate association of thought, though she was very far from conceiving it to be of equal pain.

They had no conversation together, no intercourse but what the commonest civility required. Once so much to each other! Now nothing! There \textit{had} been a time, when of all the large party now filling the drawing-room at Uppercross, they would have found it most difficult to cease to speak to one another. With the exception, perhaps, of Admiral and Mrs Croft, who seemed particularly attached and happy, (Anne could allow no other exceptions even among the married couples), there could have been no two hearts so open, no tastes so similar, no feelings so in unison, no countenances so beloved. Now they were as strangers; nay, worse than strangers, for they could never become acquainted. It was a perpetual estrangement.

When he talked, she heard the same voice, and discerned the same mind. There was a very general ignorance of all naval matters throughout the party; and he was very much questioned, and especially by the two Miss Musgroves, who seemed hardly to have any eyes but for him, as to the manner of living on board, daily regulations, food, hours, \&c., and their surprise at his accounts, at learning the degree of accommodation and arrangement which was practicable, drew from him some pleasant ridicule, which reminded Anne of the early days when she too had been ignorant, and she too had been accused of supposing sailors to be living on board without anything to eat, or any cook to dress it if there were, or any servant to wait, or any knife and fork to use.

From thus listening and thinking, she was roused by a whisper of Mrs Musgrove's who, overcome by fond regrets, could not help saying—

»Ah! Miss Anne, if it had pleased Heaven to spare my poor son, I dare say he would have been just such another by this time.«

Anne suppressed a smile, and listened kindly, while Mrs Musgrove relieved her heart a little more; and for a few minutes, therefore, could not keep pace with the conversation of the others.

When she could let her attention take its natural course again, she found the Miss Musgroves just fetching the Navy List (their own navy list, the first that had ever been at Uppercross), and sitting down together to pore over it, with the professed view of finding out the ships that Captain Wentworth had commanded.

»Your first was the Asp, I remember; we will look for the Asp.«

»You will not find her there. Quite worn out and broken up. I was the last man who commanded her. Hardly fit for service then. Reported fit for home service for a year or two, and so I was sent off to the West Indies.«

The girls looked all amazement.

»The Admiralty,« he continued, »entertain themselves now and then, with sending a few hundred men to sea, in a ship not fit to be employed. But they have a great many to provide for; and among the thousands that may just as well go to the bottom as not, it is impossible for them to distinguish the very set who may be least missed.«

»Phoo! phoo!« cried the Admiral, »what stuff these young fellows talk! Never was a better sloop than the Asp in her day. For an old built sloop, you would not see her equal. Lucky fellow to get her! He knows there must have been twenty better men than himself applying for her at the same time. Lucky fellow to get anything so soon, with no more interest than his.«

»I felt my luck, Admiral, I assure you;« replied Captain Wentworth, seriously. »I was as well satisfied with my appointment as you can desire. It was a great object with me at that time to be at sea; a very great object, I wanted to be doing something.«

»To be sure you did. What should a young fellow like you do ashore for half a year together? If a man had not a wife, he soon wants to be afloat again.«

»But, Captain Wentworth,« cried Louisa, »how vexed you must have been when you came to the Asp, to see what an old thing they had given you.«

»I knew pretty well what she was before that day;« said he, smiling. »I had no more discoveries to make than you would have as to the fashion and strength of any old pelisse, which you had seen lent about among half your acquaintance ever since you could remember, and which at last, on some very wet day, is lent to yourself. Ah! she was a dear old Asp to me. She did all that I wanted. I knew she would. I knew that we should either go to the bottom together, or that she would be the making of me; and I never had two days of foul weather all the time I was at sea in her; and after taking privateers enough to be very entertaining, I had the good luck in my passage home the next autumn, to fall in with the very French frigate I wanted. I brought her into Plymouth; and here another instance of luck. We had not been six hours in the Sound, when a gale came on, which lasted four days and nights, and which would have done for poor old Asp in half the time; our touch with the Great Nation not having much improved our condition. Four-and-twenty hours later, and I should only have been a gallant Captain Wentworth, in a small paragraph at one corner of the newspapers; and being lost in only a sloop, nobody would have thought about me.« Anne's shudderings were to herself alone; but the Miss Musgroves could be as open as they were sincere, in their exclamations of pity and horror.

»And so then, I suppose,« said Mrs Musgrove, in a low voice, as if thinking aloud, »so then he went away to the Laconia, and there he met with our poor boy. Charles, my dear,« (beckoning him to her), »do ask Captain Wentworth where it was he first met with your poor brother. I always forgot.«

»It was at Gibraltar, mother, I know. Dick had been left ill at Gibraltar, with a recommendation from his former captain to Captain Wentworth.«

»Oh! but, Charles, tell Captain Wentworth, he need not be afraid of mentioning poor Dick before me, for it would be rather a pleasure to hear him talked of by such a good friend.«

Charles, being somewhat more mindful of the probabilities of the case, only nodded in reply, and walked away.

The girls were now hunting for the Laconia; and Captain Wentworth could not deny himself the pleasure of taking the precious volume into his own hands to save them the trouble, and once more read aloud the little statement of her name and rate, and present non-commissioned class, observing over it that she too had been one of the best friends man ever had.

»Ah! those were pleasant days when I had the Laconia! How fast I made money in her. A friend of mine and I had such a lovely cruise together off the Western Islands. Poor Harville, sister! You know how much he wanted money: worse than myself. He had a wife. Excellent fellow. I shall never forget his happiness. He felt it all, so much for her sake. I wished for him again the next summer, when I had still the same luck in the Mediterranean.«

»And I am sure, Sir,« said Mrs Musgrove, »it was a lucky day for \textit{us}, when you were put captain into that ship. \textit{We} shall never forget what you did.«

Her feelings made her speak low; and Captain Wentworth, hearing only in part, and probably not having Dick Musgrove at all near his thoughts, looked rather in suspense, and as if waiting for more.

»My brother,« whispered one of the girls; »mamma is thinking of poor Richard.«

»Poor dear fellow!« continued Mrs Musgrove; »he was grown so steady, and such an excellent correspondent, while he was under your care! Ah! it would have been a happy thing, if he had never left you. I assure you, Captain Wentworth, we are very sorry he ever left you.«

There was a momentary expression in Captain Wentworth's face at this speech, a certain glance of his bright eye, and curl of his handsome mouth, which convinced Anne, that instead of sharing in Mrs Musgrove's kind wishes, as to her son, he had probably been at some pains to get rid of him; but it was too transient an indulgence of self-amusement to be detected by any who understood him less than herself; in another moment he was perfectly collected and serious, and almost instantly afterwards coming up to the sofa, on which she and Mrs Musgrove were sitting, took a place by the latter, and entered into conversation with her, in a low voice, about her son, doing it with so much sympathy and natural grace, as shewed the kindest consideration for all that was real and unabsurd in the parent's feelings.

They were actually on the same sofa, for Mrs Musgrove had most readily made room for him; they were divided only by Mrs Musgrove. It was no insignificant barrier, indeed. Mrs Musgrove was of a comfortable, substantial size, infinitely more fitted by nature to express good cheer and good humour, than tenderness and sentiment; and while the agitations of Anne's slender form, and pensive face, may be considered as very completely screened, Captain Wentworth should be allowed some credit for the self-command with which he attended to her large fat sighings over the destiny of a son, whom alive nobody had cared for.

Personal size and mental sorrow have certainly no necessary proportions. A large bulky figure has as good a right to be in deep affliction, as the most graceful set of limbs in the world. But, fair or not fair, there are unbecoming conjunctions, which reason will patronize in vain—which taste cannot tolerate—which ridicule will seize.

The Admiral, after taking two or three refreshing turns about the room with his hands behind him, being called to order by his wife, now came up to Captain Wentworth, and without any observation of what he might be interrupting, thinking only of his own thoughts, began with—

»If you had been a week later at Lisbon, last spring, Frederick, you would have been asked to give a passage to Lady Mary Grierson and her daughters.«

»Should I? I am glad I was not a week later then.«

The Admiral abused him for his want of gallantry. He defended himself; though professing that he would never willingly admit any ladies on board a ship of his, excepting for a ball, or a visit, which a few hours might comprehend.

»But, if I know myself,« said he, »this is from no want of gallantry towards them. It is rather from feeling how impossible it is, with all one's efforts, and all one's sacrifices, to make the accommodations on board such as women ought to have. There can be no want of gallantry, Admiral, in rating the claims of women to every personal comfort \textit{high}, and this is what I do. I hate to hear of women on board, or to see them on board; and no ship under my command shall ever convey a family of ladies anywhere, if I can help it.«

This brought his sister upon him.

»Oh! Frederick! But I cannot believe it of you.—All idle refinement!—Women may be as comfortable on board, as in the best house in England. I believe I have lived as much on board as most women, and I know nothing superior to the accommodations of a man-of-war. I declare I have not a comfort or an indulgence about me, even at Kellynch Hall,« (with a kind bow to Anne), »beyond what I always had in most of the ships I have lived in; and they have been five altogether.«

»Nothing to the purpose,« replied her brother. »You were living with your husband, and were the only woman on board.«

»But you, yourself, brought Mrs Harville, her sister, her cousin, and three children, round from Portsmouth to Plymouth. Where was this superfine, extraordinary sort of gallantry of yours then?«

»All merged in my friendship, Sophia. I would assist any brother officer's wife that I could, and I would bring anything of Harville's from the world's end, if he wanted it. But do not imagine that I did not feel it an evil in itself.«

»Depend upon it, they were all perfectly comfortable.«

»I might not like them the better for that perhaps. Such a number of women and children have no \textit{right} to be comfortable on board.«

»My dear Frederick, you are talking quite idly. Pray, what would become of us poor sailors' wives, who often want to be conveyed to one port or another, after our husbands, if everybody had your feelings?«

»My feelings, you see, did not prevent my taking Mrs Harville and all her family to Plymouth.«

»But I hate to hear you talking so like a fine gentleman, and as if women were all fine ladies, instead of rational creatures. We none of us expect to be in smooth water all our days.«

»Ah! my dear,« said the Admiral, »when he has got a wife, he will sing a different tune. When he is married, if we have the good luck to live to another war, we shall see him do as you and I, and a great many others, have done. We shall have him very thankful to anybody that will bring him his wife.«

»Ay, that we shall.«

»Now I have done,« cried Captain Wentworth. »When once married people begin to attack me with,—»Oh! you will think very differently, when you are married.« I can only say, »No, I shall not;« and then they say again, »Yes, you will,« and there is an end of it.«

He got up and moved away.

»What a great traveller you must have been, ma'am!« said Mrs Musgrove to Mrs Croft.

»Pretty well, ma'am, in the fifteen years of my marriage; though many women have done more. I have crossed the Atlantic four times, and have been once to the East Indies, and back again, and only once; besides being in different places about home: Cork, and Lisbon, and Gibraltar. But I never went beyond the Streights, and never was in the West Indies. We do not call Bermuda or Bahama, you know, the West Indies.«

Mrs Musgrove had not a word to say in dissent; she could not accuse herself of having ever called them anything in the whole course of her life.

»And I do assure you, ma'am,« pursued Mrs Croft, »that nothing can exceed the accommodations of a man-of-war; I speak, you know, of the higher rates. When you come to a frigate, of course, you are more confined; though any reasonable woman may be perfectly happy in one of them; and I can safely say, that the happiest part of my life has been spent on board a ship. While we were together, you know, there was nothing to be feared. Thank God! I have always been blessed with excellent health, and no climate disagrees with me. A little disordered always the first twenty-four hours of going to sea, but never knew what sickness was afterwards. The only time I ever really suffered in body or mind, the only time that I ever fancied myself unwell, or had any ideas of danger, was the winter that I passed by myself at Deal, when the Admiral (\textit{Captain} Croft then) was in the North Seas. I lived in perpetual fright at that time, and had all manner of imaginary complaints from not knowing what to do with myself, or when I should hear from him next; but as long as we could be together, nothing ever ailed me, and I never met with the smallest inconvenience.«

»Aye, to be sure. Yes, indeed, oh yes! I am quite of your opinion, Mrs Croft,« was Mrs Musgrove's hearty answer. »There is nothing so bad as a separation. I am quite of your opinion. \textit{I} know what it is, for Mr Musgrove always attends the assizes, and I am so glad when they are over, and he is safe back again.«

The evening ended with dancing. On its being proposed, Anne offered her services, as usual; and though her eyes would sometimes fill with tears as she sat at the instrument, she was extremely glad to be employed, and desired nothing in return but to be unobserved.

It was a merry, joyous party, and no one seemed in higher spirits than Captain Wentworth. She felt that he had every thing to elevate him which general attention and deference, and especially the attention of all the young women, could do. The Miss Hayters, the females of the family of cousins already mentioned, were apparently admitted to the honour of being in love with him; and as for Henrietta and Louisa, they both seemed so entirely occupied by him, that nothing but the continued appearance of the most perfect good-will between themselves could have made it credible that they were not decided rivals. If he were a little spoilt by such universal, such eager admiration, who could wonder?

These were some of the thoughts which occupied Anne, while her fingers were mechanically at work, proceeding for half an hour together, equally without error, and without consciousness. \textit{Once} she felt that he was looking at herself, observing her altered features, perhaps, trying to trace in them the ruins of the face which had once charmed him; and \textit{once} she knew that he must have spoken of her; she was hardly aware of it, till she heard the answer; but then she was sure of his having asked his partner whether Miss Elliot never danced? The answer was, »Oh, no; never; she has quite given up dancing. She had rather play. She is never tired of playing.« Once, too, he spoke to her. She had left the instrument on the dancing being over, and he had sat down to try to make out an air which he wished to give the Miss Musgroves an idea of. Unintentionally she returned to that part of the room; he saw her, and, instantly rising, said, with studied politeness—

»I beg your pardon, madam, this is your seat;« and though she immediately drew back with a decided negative, he was not to be induced to sit down again.

Anne did not wish for more of such looks and speeches. His cold politeness, his ceremonious grace, were worse than anything.